\chapter{Fixed Income Securities}
A fixed income security is an investment that provides a return in the form of fixed periodic
interest payments and the eventual return of principal at maturity. The payments of a fixed income
security are known in advance. Bonds are the most common form of fixed income securities.

A derivative is an instrument which value is derived from a primary instrument (stocks, currency,
bonds, other derivatives etc.) also called underlying instrument. (Derivative contracts divided into
OTC-derivatives (over the counter) and stock exchange derivatives and symmetric and asymmetric
derivatives). Derivatives are priced using Martingale pricing theory. Derivatives are effective
tools for controlling risk. With them you can separate risk and trade it separately. Thus, can sell
unwanted risk or buy risk to diversify the portfolio.

\section{Forward contract}
In a forward contract you agree to buy or sell an asset in a future period at a pre specified price,
all agreed on the trade date. Symmetric, i.e. when the agreement is made, the value is zero for both
parties. (Used to remove different kinds of risk (currency, interest rate), and speculating in
rising or falling prices. Often OTC-product. Synthetic product, i.e. you do not own the underlying,
hence a running compensation is made. (Most common: Currency forwards, forward rate agreements.)

\subsection{Pricing}
A forward with maturity $T$ and delivery price $K$ provides a payoff of $P_{T}-K$ at time $T$ and no
other payments. Here, $P$ is the underlying (typically the price of an asset or a specific interest
rate). Then the time $t<T$ value of the payoff is
\[
V_{t}=\mathbb{E}_{t}^{\mathbb{Q}}\left[\mathrm{e}^{-\int_{t}^{T} r_{u} \mathrm{~d} u}\left(P_{T}-K\right)\right]=\mathbb{E}_{t}^{\mathbb{Q}}\left[\mathrm{e}^{-\int_{t}^{T} r_{u} \mathrm{~d} u} P_{T}\right]-K B_{t}^{T}
\]
By definition the forward price for immediate delivery is identical to the value of the underlying,
i.e. $F_{T}^{T}=P_{T}$. The forward price is the value of $K$ that solves $V_{t}=0$, i.e. the
forward price is
\[
F_{t}^{T}=\frac{\mathbb{E}_{t}^{\mathbb{Q}}\left[\mathrm{e}^{-\int_{t}^{T} r_{u} \mathrm{~d} u} P_{T}\right]}{B_{t}^{T}}
\]
If $P$ is the price of a traded asset that pays only at time $T,
F_{t}^{T}=\frac{P_{t}}{B_{t}^{T}}$.

\section{European options}
European options provide a right to trade the underlying asset on expiration date for a pre
specified (at time $t$ ) price $K$ (American options can be exercised on any date). For this
opportunity a premium is paid. (Call is right to buy, Put is right to sell. Asymmetric due to traded
premium on agreement. With a call option you speculate on rising prices, put is falling prices e.g.
Can be used to remove risk for future stock trades.)

\subsection{Pricing}
Consider a call option with expiration date $T$ and strike $K$ on the underlying $P_{t}$. The payoff
at time $T$ (and hence the call price at $T$ ) is
\[
C_{T}=\max \left\{P_{T}-K, 0\right\}
\]
Remember, there exists an equivalent martingale measure $\mathbb{Q}$ so the asset price $P_{T}$ is a
martingale (under the defined probability measure) using the bank account $A_{t}= \exp
\left(\int_{0}^{t} r_{s} \mathrm{~d} s\right)$ as numeraire, i.e.
\[
\frac{P_{t}}{A_{t}}=\mathbb{E}_{t}^{\mathbb{Q}}\left[\frac{P_{T}}{A_{T}}\right]
\]
Using $B_{t}^{T}$ as numeraire (recall $B_{T}^{T}=1$ ) we have
\[
P_{t}=B_{t}^{T} \mathbb{E}_{t}^{\mathbb{Q}^{T}}\left[P_{T}\right]
\]
with $\mathbb{Q}^{T}$ being the $T$-forward MG measure. Using this measure we can price the call
option by
\[
\begin{aligned}
C_{t} & =B_{t}^{T} \mathbb{E}_{t}^{\mathbb{Q}^{T}}\left[\max \left\{P_{T}-K, 0\right\}\right] \\
& =B_{t}^{T} \mathbb{E}_{t}^{\mathbb{Q}^{T}}\left[\1_{\left\{P_{T}>K\right\}}\left(P_{T}-K\right)\right] \\
& =B_{t}^{T}\left(\mathbb{E}_{t}^{\mathbb{Q}^{T}}\left[\1_{\left\{P_{T}>K\right\}} P_{T}\right]-K \mathbb{Q}^{T}\left[P_{T}>K\right]\right) \\
& =B_{t}^{T} \mathbb{E}_{t}^{\mathbb{Q}^{T}}\left[\1_{\left\{P_{T}>K\right\}} P_{T}\right]-K B_{t}^{T} \mathbb{Q}^{T}\left[P_{T}>K\right] \\
& =P_{t} \mathbb{E}_{t}^{\mathbb{Q}^{P}}\left[\1_{\left\{P_{T}>K\right\}}\right]-K B_{t}^{T} \mathbb{Q}_{t}^{T}\left[P_{T}>K\right] \\
& =P_{t} \mathbb{Q}_{t}^{P}\left[P_{T}>K\right]-K B_{t}^{T} \mathbb{Q}_{t}^{T}\left[P_{T}>K\right]
\end{aligned}
\]
where $\mathbb{Q}^{P}$ is the MG measure using $P_{t}$ as numeraire.

\section{Additional stuff}
\begin{description}
\item[Future] 
\end{description}
When buying a future you get compensated everyday (continously) for the difference in the
new future price. Therefore, you get implicitly the continous rate through the future price. (With a
forward you only take into account the expected future rate. Thus, the difference is only the
discounted covariance between the zero-coupon bond and the payment price at time $T$.) Future price:
$\phi_{t}^{T}=\mathbb{E}^{\mathbb{Q}}\left[P_{T}\right]$.i

\begin{description}
\item[Caps] 
\end{description}
A cap protects a floating-rate borrower against the risk of paying very high interest rates.
Work as a ceiling on the interest rate, compensating the holder if the rate rises above a specified
level. Consists of a series of underlying options called caplets (call option on the spot interest
rate).

\begin{description}
\item[Floors] 
\end{description}
A floor protects a floating-rate lender against the risk of receiving very low interest
rates. Works like caps, consists of floorlets (put option on floating interest rate with delayed
payment of the payoff).

\begin{description}
\item[Collars] 
\end{description}
A combination of caps and floors. Ensures that the interest rate payments on a floating
rate borrowing arrangement stay between two prespecified levels. (Can be seen as a portfolio of a
long position in a cap with rate $K_{c}$ and a short position in floor with rate $K_{f}$.)

\begin{description}
\item[Swaps] 
\end{description}
An interest rate swap is an exchange of two cash flow streams determined by certain interest
rates. A plain vanilla interest rate swap takes two forms:

\begin{itemize}
  \item Payer swap: Pays a stream of fixed rate payments and recieves a stream of floating rate
    payments.
  \item Receiver swap: Pays a stream of floating rate payments and receives a stream of fixed rate
    payments.
\end{itemize}
Forward swap: Agreement to enter into a swap with a future starting date and fixed rate already
set.

\begin{description}
\item[Swaptions] 
\end{description}
A European swaption gives its holder the right at the expiry date to enter into an
interest rate swap that has a given fixed rate. A payer swaption gives the right to enter into a
payer swap, and likewise for receiver swaption.

